% !TeX root = ../papers/nc-p3.tex
\section{六次の例外ファイバーと九次で必要になる中心の選び直し}
\label{sec:rank-fifteen-fixed-center-obstruction}
\begin{proposition}
標数3の非分岐幾何単元対では、前節の第1,2,4,5列が四始点とも
六次まで所定の次元を与える表示を常に選べる。
しかしこの同じ四列では、非分岐域のある表示で九次の商の次元が
全始点で23となる。その表示を保ったまま第4列を第3列に替えると、
四始点とも次数0から13までの商の次元は
\[
1,4,10,20,31,40,44,40,31,20,10,4,1,0
\]
となる。また同じ表示に、九次まで所定の次元を持ち、十次で初めて増大する別の四中心もある。
\end{proposition}
\begin{proof}
前節の六次の終結式で残る横座標を $X^2+X+2=0$ とする。
$G$ の十二基底係数を $\mathbb F_9[B]$ に特殊化したモニック最大公約元は
$B-(X+1)$ である。前節の検証コードは同じ $G$ からこの多項式を計算する。
従って指定四列の一般主元が全十二表示で同時に消えうるのはこの共役な点対だけである。
ここでいう消失は、四始点の主元の積の消失を表す。

$a^4+2a^3+2=0$ とした $\mathbb F_{81}$ で、次の一表示を取る。
\[
 X=a^3+a^2+2,\qquad c=a^2+2,\qquad v=a^3+a^2+a.
\]
実際に $c^2=X$, $(v^3-v)^2=X+1$ であり、$K=1/v$, $H=c(K-v)$
で六根の二乗の積は非零である。共役点はFrobeniusで得る。
\path{computations/verify_rank_fifteen_fixed_center_obstruction.py} は普遍的な
五自由列の語係数をこの $K,H$ に特殊化し、体上で主元を選び直して
前節の末尾語の再帰を実行する。第1,2,4,5列では四始点とも八次まで上の次元に一致する。
九次の文字別次元は、始点0,2で $(6,6,6,5)$、始点1,3で $(5,6,6,6)$ である。
これで残っていた点対の六次と九次の増大が同時に分かる。
第1,2,3,5列で同じ再帰を十三次まで行うと表示した列になる。
さらに五列を $z_1,\ldots,z_5$ と書き、$z_1+az_3,z_2+a^2z_3,z_4+a^3z_3,z_5+a^4z_3$ を取る。
この四中心は独立である。全始点で九次まで期待値に一致し、十次から十三次の各次元は16となる。
特に九次までの一致だけでは、選んだ四中心が所定の商を与えるかは決まらない。
\end{proof}
